\documentclass{article}
\usepackage[T1]{fontenc}
\usepackage{lmodern}
\usepackage{graphicx}
\usepackage{amsmath,amssymb}
\usepackage[a4paper,margin=2cm]{geometry}
\usepackage[absolute,overlay]{textpos}
\setlength{\TPHorizModule}{1mm}
\setlength{\TPVertModule}{1mm}
\usepackage[indonesian]{babel}
\usepackage{hyperref}
\setlength{\emergencystretch}{2em}
% unit: Halaman 6

\begin{document}

% === Halaman 6 (python/native) ===

4.  <Zn, > : grup dengan himpunan integer modulo n, yaitu Zn = \{0, 1, 2, …, n – 1\} 

dan  adalah operasi penjumlahan modulo n. 

    Contoh: n = 5, Zn = \{0, 1, 2, 3, 4\},   (3  4) =  (3 + 4) mod 5 = 2

  <Zp, > : grup dengan himpunan integer modulo p, p adalah bilangan prima, 

yaitu Zp = \{0, 1, 2, …, p – 1\} dan   adalah operasi penjumlahan modulo p. 

 <Z*p,  > : dengan himpunan integer bukan nol, p adalah bilangan prima, yaitu 

Z*p = \{1, 2, …, p – 1\} dan  adalah operasi perkalian modulo p. 

6

Rinaldi Munir/II4021 Kriptografi

\end{document}
